\documentclass[a4paper,12pt]{article}
\usepackage[utf8]{inputenc}
\usepackage[T1]{fontenc}
\usepackage[ngerman]{babel}
\usepackage{geometry}
\geometry{margin=2.5cm}
\usepackage{amsmath}
\usepackage{hyperref}
\hypersetup{colorlinks=true, linkcolor=blue, urlcolor=blue, citecolor=blue}

\begin{document}

\section*{Neural Operator: Learning Maps Between Function Spaces with Applications to PDEs}

\noindent\textbf{Autoren:} Nikola Kovachki, Zongyi Li, Burigede Liu,
Kamyar Azizzadenesheli, Kaushik Bhattacharya, Andrew Stuart, Anima Anandkumar\\[2pt]
\textbf{Erschienen in:} Journal of Machine Learning Research, Vol.\,24 (2023), S.\,1--97\\
\textbf{Eingereicht:} Dezember 2021 \quad \textbf{Veröffentlicht:} Dezember 2022

% -----------------------------------------------------------------
\subsection*{Kernbeitrag}

Diese Arbeit schlägt ein allgemeines Deep-Learning-Framework vor, das klassische
neuronale Netze von Abbildungen zwischen endlichdimensionalen euklidischen Räumen
auf Abbildungen zwischen \emph{unendlichdimensionalen Funktionenräumen} verallgemeinert.
Die resultierende Modellklasse wird als \textbf{Neuronaler Operator} bezeichnet.
Motiviert wird dieser Ansatz durch PDE-Anwendungen, bei denen die Eingangs- und
Ausgangsgröße des zu erlernenden Lösungsoperators Funktionen und keine Vektoren sind
-- zum Beispiel die Abbildung von Anfangs- oder Randbedingungen auf die zugehörige
Lösungsfunktion.

% -----------------------------------------------------------------
\subsection*{Theoretische Grundlagen}

\paragraph{Diskretisierungsinvarianz.}
Ein zentrales Konzept der Arbeit ist die \emph{Diskretisierungsinvarianz}:
Ein Modell mit fester Parameteranzahl gilt als diskretisierungsinvariant, wenn es
(i) beliebige Diskretisierungen der Eingangsfunktion akzeptiert,
(ii) die Ausgabefunktion an beliebigen Punkten auswerten kann, und
(iii) im Grenzfall immer feinerer Diskretisierung gegen einen kontinuierlichen
Operator konvergiert.
Die Autoren zeigen formal, dass klassische CNNs und Standard-Neuronalnetze diese
Eigenschaft \emph{nicht} besitzen -- ihr Fehler wächst bei Gitterverfeinerung --,
während neuronale Operatoren sie garantieren.
Für die Masterarbeit bedeutet das, dass ein FNO-Modell prinzipiell auf einer
niedrigeren Trainingsauflösung trainiert und dann auf $256{\times}256$-Gittern
ausgewertet werden kann.

\paragraph{Universelle Approximation.}
Theorem\,11 der Arbeit beweist, dass neuronale Operatoren jeden stetigen Operator
$G^\dagger: \mathcal{A} \to \mathcal{U}$ zwischen Banachräumen gleichmäßig auf
kompakten Mengen approximieren können. Dies wird sowohl für $L^p$- und
Sobolev-Räume als auch für Räume stetiger Funktionen gezeigt und stellt die
einzige bekannte Modellklasse dar, die gleichzeitig Diskretisierungsinvarianz
\emph{und} universelle Approximation garantiert.

% -----------------------------------------------------------------
\subsection*{Architektur}

Der neuronale Operator der Arbeit folgt der allgemeinen Struktur
\[
  G_\theta = Q \circ \sigma_T(W_{T-1} + K_{T-1} + b_{T-1})
              \circ \cdots \circ
              \sigma_1(W_0 + K_0 + b_0) \circ P,
\]
wobei $P$ einen punktweisen Lifting-Schritt, $K_t$ einen nichtlokalen
Integral-Kern-Operator, $W_t$ einen lokalen linearen Operator und $Q$
eine abschließende Projektion bezeichnet.
Insgesamt werden vier konkrete Parametrisierungen des Kernoperators vorgestellt:

\begin{itemize}
  \item \textbf{GNO} (Graph Neural Operator): Approximiert das Integral über
        Nystrom-Stichproben auf einem Graphen; flexibel für unstrukturierte Gitter.
  \item \textbf{LNO} (Low-rank Neural Operator): Dekomponiert den Kern in eine
        Produkt-Form mit linearer Komplexität $\mathcal{O}(J)$.
  \item \textbf{MGNO} (Multipole Graph Neural Operator): Kombiniert lokale
        Graphstruktur mit hierarchischer Mehrskalendekomposition à la Fast
        Multipole Method; ebenfalls $\mathcal{O}(J)$.
  \item \textbf{FNO} (Fourier Neural Operator): Parametrisiert den Kern direkt
        im Frequenzraum über die FFT mit quasilinearer Komplexität
        $\mathcal{O}(J \log J)$; derzeit bester Ansatz auf uniformen Gittern.
\end{itemize}

Jede FNO-Schicht führt eine Fouriertransformation durch, multipliziert die
niedrigfrequenten Moden mit einer lernbaren Matrix $R \in \mathbb{C}^{k_\text{max}
\times m \times n}$ und transformiert zurück, addiert dabei einen lokalen
residualen Pfad $W v(x)$:
\[
  u(x) = \mathcal{F}^{-1}\!\bigl(R \cdot \mathcal{F}(v)\bigr)(x) + W v(x) + b(x).
\]

% -----------------------------------------------------------------
\subsection*{Numerische Ergebnisse}

Die Arbeit benchmarkt alle vier Methoden auf vier Testproblemen: der
1D-Poisson-Gleichung, Darcy-Strömung, Burgers-Gleichung und Navier-Stokes.
Wichtige Befunde für die Masterarbeit:

\begin{itemize}
  \item Auf der \textbf{2D-Darcy-Strömung} erreicht FNO einen relativen $L_2$-Fehler
        von $1{,}08\,\%$ (Auflösung $s=85$), während U-Net nicht direkt verglichen
        wird, aber FCN auf größeren Gittern deutlich schlechter abschneidet.
  \item Auf der \textbf{Navier-Stokes-Gleichung} ($\nu = 10^{-3}$, $N=1000$,
        $T=50$) erzielt FNO-3D einen relativen Fehler von $0{,}86\,\%$, während
        \textbf{U-Net denselben Fehler von $2{,}45\,\%$} aufweist -- also etwa
        dreimal höher. Dieser direkte Vergleich ist für den Architekturvergleich
        in der Masterarbeit unmittelbar relevant.
  \item FNO ist gegenüber U-Net beim \textbf{Zero-Shot Super-Resolution}-Transfer
        (Training auf $64{\times}64$, Inferenz auf $256{\times}256$) der einzige
        Ansatz, der keine Fehlerdegradierung zeigt.
  \item Im \textbf{Bayes'schen Inversionsproblem} ist FNO als Surrogatmodell
        ca.\ 440-mal schneller als der traditionelle Löser (0,005\,s vs.\ 2,2\,s
        pro Evaluierung), was MCMC-Sampling erst praktikabel macht.
\end{itemize}

% -----------------------------------------------------------------
\subsection*{Beziehung zu anderen Methoden}

Die Autoren zeigen außerdem, dass DeepONets ein Spezialfall des
neuronalen Operators sind, bei dem die Diskretisierungsinvarianz gebrochen wird
(Proposition~5). Transformers mit Self-Attention lassen sich als kontinuierliche
Grenzfälle des nichtlinearen Kernoperators interpretieren (Proposition~6).
CNNs approximieren im Kontinuumslimit Differentialoperatoren, verlieren aber
bei Gitterverfeinerung die Konsistenz -- ein wichtiger Unterschied zu FNO.

% -----------------------------------------------------------------
\subsection*{Bezug zur Masterarbeit}

\paragraph{Theoretische Einordnung des FNO.}
Die Masterarbeit vergleicht UNet- und FNO-Architekturen als Surrogatmodelle für
LaMEM-Simulationen von Stokes-Strömungsfeldern um sedimentierende Kristalle.
Kovachki et al.\ liefern die formale Grundlage dafür: Der FNO approximiert den
Lösungsoperator $G^\dagger: a \mapsto (v_x, v_y)$ -- also die Abbildung von
Kristallgeometrie und -position auf die Geschwindigkeitsfelder -- als neuronalen
Operator im Sinne von Theorem\,11. Die universelle Approximationseigenschaft
begründet, warum ein FNO prinzipiell in der Lage ist, die nichtlokalen
Druckwechselwirkungen zwischen mehreren Kristallen zu erlernen, auch wenn während
des Trainings nur $N$ Kristalle vorhanden waren.

\paragraph{Diskretisierungsinvarianz und Generalisierung.}
Die zentrale Forschungsfrage der Masterarbeit -- ob ein auf $N$ Kristallen
trainiertes Netz auf Konfigurationen mit $N{+}m$ Kristallen generalisiert --
ist eng mit der Diskretisierungsinvarianz verknüpft: Da der FNO den
Funktionsraum und nicht eine spezifische Diskretisierung erlernt, sollte er
robuster gegenüber veränderten Eingabestrukturen sein als ein klassisches UNet,
das fest an das Trainingsgitter gebunden ist. Diese Eigenschaft des FNO stützt
die Hypothese, dass er für das $N$-zu-$(N{+}m)$-Generalisierungsproblem
geeigneter sein könnte.

\paragraph{FNO vs.\ U-Net für Strömungsfelder.}
Die in Tabelle~4 der Arbeit gezeigte Überlegenheit des FNO gegenüber U-Net
auf der Navier-Stokes-Gleichung (Fehlerreduktion um Faktor ${\sim}3$) liefert
eine empirische Baseline-Erwartung: Falls ähnliche Muster auf dem Stokes-Problem
der Masterarbeit auftreten, wäre dies ein starker Hinweis auf die strukturelle
Eignung von FNO für Strömungssurrogate. Gleichzeitig ist zu beachten, dass das
Stokes-Regime ($Re \ll 1$) weit glattere Lösungen erzeugt als turbulente
Navier-Stokes-Strömungen, was dem UNet durch seine Multiskalen-Encoder-Decoder-
Struktur zugutekommen könnte.

\paragraph{Residuales Lernen und physikalische Konsistenz.}
Der in dieser Arbeit verwendete additive residuale Pfad $Wv(x)$ in jeder FNO-Schicht
ist konzeptuell verwandt mit dem in der Masterarbeit implementierten
\emph{residualen Lernansatz} $v_\text{total} = v_\text{Stokes} + \Delta v_\text{NN}$,
bei dem das Netz lediglich die Korrektur zur analytischen Stokes-Lösung erlernt.
Beide Mechanismen nutzen strukturelles Vorwissen, um den Lernprozess zu stabilisieren.

\paragraph{Spektrale Eigenschaften.}
Die in der Arbeit diskutierten spektralen Zerfallsanalysen (Abschnitt\,7.3.2)
zeigen, dass FNO auch Hochfrequenzanteile der Lösung rekonstruieren kann,
obwohl der Kern nur auf $k_\text{max}$ Moden beschränkt ist. Dies ist für die
Masterarbeit relevant, da Strömungsfelder um Kristalle im Nahfeld scharfe
Gradienten aufweisen, die präzise Hochfrequenzinformation erfordern.

\end{document}